\section{The five-dimensional frontier}
\label{app:dimension-five}

The body studies explicit descendants in dimensions \(19\), \(38\), and
\(110\).  The opposite optimization problem begins in dimension five:
cubic-homogeneous Keller maps in dimensions three and four are known to be
invertible, so five is the first dimension in which a cubic-homogeneous
counterexample could occur.  This appendix records exact restrictions and a
surviving line-level family.  Nothing here constructs a five-dimensional
counterexample.

\subsection{Collision pencils}

Let
\[
F(x)=x+H(x),\qquad H(x)=T(x,x,x),
\]
where \(T\colon V^3\to V\) is symmetric and \(\dim V=5\).

\begin{lemma}[Collision vectors are independent]
\label{lem:n5-collision-independent}
If \(u\ne0\) and \(F(v+u)=F(v)\), then \(u\) and \(v\) are linearly
independent.
\end{lemma}

\begin{proof}
Suppose \(v=\lambda u\).  Cubic homogeneity gives
\[
 H(v+u)-H(v)=cH(u),\qquad
 c=(\lambda+1)^3-\lambda^3.
\]
The collision equation is \(u=-cH(u)\).  If \(c=0\), this forces \(u=0\),
contrary to the hypothesis.  Otherwise Euler's identity gives
\[
 JH(u)u=3H(u)=-\frac3c u,
\]
so \(JH(u)\) has a nonzero eigenvalue.  On the other hand, the Keller
identity and \(JH(tu)=t^2JH(u)\) give
\[
 \det(I+t^2JH(u))=1
\]
for every scalar \(t\).  Hence every eigenvalue of \(JH(u)\) is zero, a
contradiction.
\end{proof}

Suppose now
\[
F(v+u)=F(v)
\]
with \(u\ne0\); linear independence follows from
\cref{lem:n5-collision-independent}.  On
\(\PP(\langle u,v\rangle)\), write
\[
M(s,t)=JH(su+tv)=s^2A+stB+t^2C,
\]
where
\[
A=JH(u),\qquad B=6T(u,v,-),\qquad C=JH(v).
\]
Symmetry of \(T\) and the collision give
\begin{equation}
\label{eq:n5-plane-equations}
Bu=2Av,\qquad Bv=2Cu,\qquad
u+\frac13Au+Av+Cu=0.
\end{equation}
Conversely, the first two identities are sufficient to extend the plane
data to a symmetric trilinear tensor.  They do not ensure that the global
Jacobian is nilpotent away from the plane.

\begin{proposition}[Three generic Jordan strata]
\label{prop:n5-jordan-types}
For a five-dimensional cubic-homogeneous counterexample, the generic
nilpotent Jordan type of \(JH\) must be one of
\[
(5),\qquad(4,1),\qquad(3,2).
\]
\end{proposition}

\begin{proof}
Known low-rank results exclude generic rank at most two in a
five-dimensional counterexample, while nilpotence gives rank at most four.
The nilpotent partitions of five having rank three or four are precisely the
three displayed partitions.
\end{proof}

\subsection{A corrected sparse exclusion}

Consider the square-zero two-matrix family
\[
H_{C,Q}(z)=C\left((Qz)^{*3}-Qz^{*3}\right).
\]
After a signed-permutation normalization, let \(Q=I+R\), where \(R\) has
exactly three active rows and exactly one nonzero off-diagonal entry in each
active row.  Impose the coordinate collision
\[
(I+H_{C,Q})(e_0)=(I+H_{C,Q})(e_1).
\]

\begin{theorem}[Three-row sparse exclusion]
\label{thm:n5-sparse}
No member of this coordinate-collision ansatz whose three directed support
edges collectively use all five vertices is Keller.
\end{theorem}

\begin{proof}
For one edge \(i\to j\) with coefficient \(r\), its cubic contribution is
\[
U=(z_i+rz_j)^3-z_i^3-rz_j^3.
\]
If \(b\) is the corresponding column of the compressed coefficient matrix,
the \(z_i^2\) and \(z_iz_j\) coefficients in its contribution to the trace
are
\[
3rb_j,\qquad 6rb_i+6r^2b_j.
\]
Thus the trace equation forces \(b_i=b_j=0\).

There are \(180\) three-edge patterns with distinct source rows.
The collision and trace constraints leave \(30\), forming three orbits under
the evident vertex symmetries:
\[
\begin{array}{c|c}
A&(0\to2,\ 3\to0,\ 4\to1),\\
B&(2\to0,\ 3\to0,\ 4\to1),\\
C&(2\to0,\ 3\to1,\ 4\to2).
\end{array}
\]
Writing the edge coefficients as \(\alpha,\beta,\gamma\), exact
elimination of the collision and trace equations leaves, respectively, the
following nonzero coefficients in
\(\operatorname{tr}(JH/3)^2\):
\[
\frac{18}{\beta^2-1}\quad\text{in orbit \(A\)},\qquad
18\quad\text{in orbit \(C\)}.
\]
In the generic branch of orbit \(B\), two coefficients force both
\[
B_{1,1}=0,\qquad(\beta^3-\beta)B_{1,1}=1.
\]
The branches \(\alpha=\pm1\) and \(\beta=\pm1\) instead leave
\[
\frac{36\beta}{\beta^2-1}
\quad\text{or}\quad
\frac{36\alpha}{\alpha^2-1}.
\]
Every branch is contradictory.
\end{proof}

\begin{remark}[Scope correction]
\label{rem:n5-scope-correction}
The theorem is not the blanket statement that every support-\(\le4\)
five-dimensional ansatz fails.  Its coordinate-collision, three-row, and
one-edge-per-row hypotheses are essential.  An earlier finite-field
experiment using only coefficients \(0,\pm1\) was vacuous for this collision
and is not evidence for the theorem; the exact characteristic-zero argument
above is the proof.
\end{remark}

\subsection{Everywhere-regular Jordan type
  \texorpdfstring{\((5)\)}{(5)}}

Assume \(M(s,t)\) has Jordan type \((5)\) at every point of \(\PP^1\).
Regard it as
\[
M\colon\mathcal O_{\PP^1}^5\longrightarrow
\mathcal O_{\PP^1}^5(2).
\]
For \(K_i=\ker M^i\) and \(L_i=K_i/K_{i-1}\), the induced maps
\[
L_i\longrightarrow L_{i-1}(2)
\]
are isomorphisms.  Since the degrees sum to zero,
\begin{equation}
\label{eq:n5-filtration-degrees}
\bigl(\deg L_1,\ldots,\deg L_5\bigr)=(-4,-2,0,2,4).
\end{equation}
In particular, the right kernel is an
\(\mathcal O(-4)\)-subbundle of \(\mathcal O^5\).

If \(r\) is the coordinate-span dimension of a quartic kernel generator,
the quadratic-syzygy calculation leaves only
\[
\begin{array}{c|c}
r&\text{possible splitting}\\ \hline
3&(2,2)\\
4&(1,1,2)\\
5&(1,1,1,1),
\end{array}
\]
while \(r=2\) is impossible.

In the full-span chart, normalize the kernel to the rational normal quartic
\[
k(s,t)=(s^4,s^3t,s^2t^2,st^3,t^4)^T.
\]
Its linear syzygy matrix is
\[
K(s,t)=
\begin{pmatrix}
-t&s&0&0&0\\
0&-t&s&0&0\\
0&0&-t&s&0\\
0&0&0&-t&s
\end{pmatrix}.
\]
The linear lifting system for the regular nilpotent quotient has rank
\(48\) in \(56\) unknowns and hence an eight-dimensional solution space.
On the open set where the resulting constant \(4\times4\) matrix \(T\) is
invertible, the pencil factors as
\[
M(s,t)=R(s,t)T^{-1}K(s,t)
\]
and satisfies \(M^5=0\).

\begin{theorem}[A smooth characteristic-zero collision-line family]
\label{thm:n5-regular-family}
The everywhere-regular type-\((5)\) equations
\eqref{eq:n5-plane-equations} have a smooth one-dimensional family in
characteristic zero.
\end{theorem}

\begin{proof}[Exact computer-assisted existence proof]
In the eight-parameter full-kernel chart, the following point over
\(\F_{11}\) satisfies the integrability and collision equations:
\[
a=(8,7,1,7,2,9,0,1),
\]
\[
u=(7,6,3,10,4),\qquad
v=(8,9,7,9,1).
\]
At this point \(\det T=1\).  The gcd of the \(4\times4\) minors of
\(M(s,t)\) is one in \(\F_{11}[s,t]\), so the pencil is geometrically
regular on \(\PP^1_{\overline{\F}_{11}}\).

After fixing \(a_7=1\) and one coordinate of \(v\), the ten integrability
and five collision equations are functions of \(16\) variables.  Their
Jacobian has rank \(15\); a specified \(15\times15\) minor has determinant
\(3\) modulo \(11\).  The multivariate Hensel lemma therefore gives a
one-dimensional smooth characteristic-zero lift.  The two accompanying
dimension-five supplement directories contain independent exact Hensel
replay scripts.  They check the lift through modulus
\[
11^{12}=3138428376721.
\]
Geometric regularity is open, so it persists on the lifted family.
\end{proof}

\subsection{The first-normal obstruction}

The smooth family in \cref{thm:n5-regular-family} is a family of
collision-line data.  It does not extend to a global nilpotent Jacobian.
Choose coordinates with \(u=e_0\), \(v=e_1\), and let
\(n_0,n_1,n_2\) span the normal directions.  For each \(r\), put
\[
N_r(s,t)=
\left.\frac{d}{d\epsilon}
JH(su+tv+\epsilon n_r)\right|_{\epsilon=0}
=sP_r+tQ_r.
\]
Tensor symmetry fixes
\[
P_ru=2An_r,\quad P_rv=Bn_r,\qquad
Q_ru=Bn_r,\quad Q_rv=2Cn_r.
\]
The remaining first-normal tensor entries form a \(60\)-dimensional vector
space.  Nilpotence through first normal order requires
\begin{equation}
\label{eq:first-normal-traces}
\operatorname{tr}\bigl(M(s,t)^{k-1}N_r(s,t)\bigr)=0,
\qquad r=0,1,2,\quad k=1,\ldots,5.
\end{equation}

\begin{theorem}[First-normal obstruction on the smooth residue disk]
\label{thm:n5-first-normal}
The complete smooth \(\mathbb Z_{11}\)-residue disk through the point in
\cref{thm:n5-regular-family} has no extension through first normal order to
a symmetric cubic tensor with nilpotent Jacobian.

In collision-adapted coordinates, the equations for \(k\le4\) in
\eqref{eq:first-normal-traces} have a \(60\times60\) coefficient matrix
with determinant \(3\pmod {11}\).  Adding \(k=5\) gives
\[
\operatorname{rank}L=60,\qquad
\operatorname{rank}[L\mid b]=61,
\]
with a selected augmented \(61\times61\) minor of determinant
\(1\pmod {11}\).
\end{theorem}

\begin{proof}[Exact computer-assisted proof]
At the displayed residue point, the transformed line matrices are
\[
A=\begin{pmatrix}
6&4&0&10&2\\6&1&0&9&10\\8&0&0&6&2\\
3&8&0&5&2\\3&6&0&1&10
\end{pmatrix},
\]
\[
B=\begin{pmatrix}
8&8&1&4&7\\2&5&2&6&9\\0&2&5&4&4\\
5&4&6&9&7\\1&8&10&3&6
\end{pmatrix},
\qquad
C=\begin{pmatrix}
4&0&5&1&3\\8&1&6&6&6\\1&8&5&6&10\\
2&6&0&0&2\\4&7&9&1&1
\end{pmatrix}.
\]
The first four trace identities form a square system in the \(60\)
first-normal coefficients and have determinant \(3\).  Substituting its
unique solution into the fifth identities leaves the coefficient matrix
\[
\begin{pmatrix}
5&8&9&1&9&2&2&1&3&1\\
2&2&6&2&4&6&1&7&2&5\\
9&5&9&8&5&3&4&6&3&8
\end{pmatrix},
\]
whose first three columns have determinant \(2\).  This proves
inconsistency.  The determinants \(3\) and \(1\) are \(11\)-adic units, so
the contradiction persists throughout the residue disk.

An independent coordinate-free calculation starts with all
\[
5\binom73=175
\]
coefficients of a symmetric vector-valued cubic.  The line restrictions
have rank \(65\); adjoining the identities through \(k=4\) gives
coefficient and augmented ranks \((125,125)\); adjoining \(k=5\) gives
\((125,126)\).  The remaining \(50\) variables after \(k\le4\) are exactly
the pure-normal coefficients \(5\binom53\), which cannot affect first
normal order.  This confirms that the obstruction is independent of the
chosen normal complement.
\end{proof}

The same exact test was applied exhaustively to the rational points of the
eight-parameter full-kernel chart:
\[
\begin{array}{c|rr}
&\F_7&\F_{11}\\ \hline
\text{projective parameter points}&960800&21435888\\
\text{invertible chart matrix}&806344&19324668\\
\text{regular collision-line solutions}&6&11\\
\text{first-normal extensions}&0&0.
\end{array}
\]
Every surviving pencil is geometrically regular over the algebraic closure,
and every first-normal system has ranks \((60,61)\).  These enumerations
exclude the rational points in the displayed chart and their
good-reduction local branches; they do not classify points over finite
extensions or other kernel-span strata.

There is also an exact characteristic-zero presentation of the displayed
collision chart.  After normalizing \(a_7=v_4=1\) and clearing the matrix
denominator, the two integrability identities and the collision identity
give fifteen primitive quintics in sixteen variables.  The relevant open
condition is
\[
 \det(T)(u_3-u_4v_3)\ne0.
\]
At the recorded \(\F_{11}\)-point all fifteen equations vanish and their
Jacobian has rank fifteen, giving a smooth one-dimensional local collision
locus on this chart.  The exact polynomial list and a Macaulay2 saturation
input are included in the computational supplement.  No global saturation
or irreducible-component calculation is asserted here.

\begin{question}[Global first-normal obstruction]
\label{q:n5-global-first-normal}
Let \(C_4\) and \(\widetilde C_4\) be the coefficient and augmented matrices
of the first-normal identities through \(k=4\), and let \(C_5\) and
\(\widetilde C_5\) include the fifth identity.  Is
\[
 \operatorname{rank}\widetilde C_5>
 \operatorname{rank}C_5
\]
at every point of the regular full-kernel collision-line curve?  On the open
locus \(\det C_4\ne0\), solving the first four identities makes this the
nonvanishing of the familiar fifth-residual covariant
\[
 \Omega\in N_P^*\otimes H^0(\PP^1,\mathcal O(9)).
\]
At fifteen of the seventeen enumerated \(\F_7\)- and \(\F_{11}\)-points,
\(C_4\) is invertible and the residual map has rank three.  At the other two,
\((\operatorname{rank}C_4,\operatorname{rank}\widetilde C_4)=(59,60)\), so
incompatibility occurs already through \(k=4\).  Thus \(\Omega\) is only an
open-locus presentation; the global object is the cokernel or, equivalently,
the augmented-rank condition.  A positive answer would exclude the entire
full-kernel regular type-\((5)\) stratum, rather than only the known residue
disk and rational points.
\end{question}

\subsection{A global rank-drop obstruction}

Let \(M(x)\) be a homogeneous quadratic nilpotent \(5\times5\) matrix on
\(\PP^4\), generically of type \((5)\), and put
\[
D=\set{[x]\in\PP^4:\operatorname{rank}M(x)\le3}.
\]

\begin{theorem}[Codimension-two rank drop]
\label{thm:n5-rank-drop}
One has
\[
\operatorname{codim}_{\PP^4}D\le2.
\]
If \(M\) extends an everywhere-regular collision line, then
\[
\operatorname{codim}_{\PP^4}D=2
\]
and \(D\) is disjoint from that line.
\end{theorem}

\begin{proof}
Suppose \(\operatorname{codim}D\ge3\) and put
\(U=\PP^4\setminus D\).  The regular kernel filtration on \(U\) has line
quotients
\[
L_i\simeq\mathcal O_U(-4+2(i-1))
\]
by the same calculation as \eqref{eq:n5-filtration-degrees}.  Removing a
codimension-at-least-three set preserves \(\operatorname{Pic}\) and
\(\operatorname{CH}^2\).  Multiplicativity of Chern classes in the
filtration of \(\mathcal O_U^5\) would give
\[
0=c_2(\mathcal O_U^5)
=\sum_{0\le i<j\le4}(-4+2i)(-4+2j)h^2
=-20h^2,
\]
contradicting \(h^2\ne0\) in \(\operatorname{CH}^2(U)\).

An everywhere-regular collision line is disjoint from \(D\).  A divisor in
\(\PP^4\) cannot be disjoint from a projective line, so \(D\) cannot have
codimension one.  The preceding bound then forces codimension two.
\end{proof}

\begin{remark}[The remaining problem]
\Cref{thm:n5-first-normal} eliminates the explicit smooth full-kernel
residue disk, but not every regular type-\((5)\) collision-line component.
The kernel-span-three and kernel-span-four strata and the generic Jordan
types \((4,1)\) and \((3,2)\) also remain open.  Every successful global
extension must acquire the codimension-two rank-drop surface of
\cref{thm:n5-rank-drop}.  Thus dimension five remains possible, but the
known regular line family is not the restriction of a counterexample.
\end{remark}
